\begin{proof}

\begin{enumerate}[label={$\langle1\rangle$\arabic*.},leftmargin=*]
\item Let $D_1,\ldots,D_r$ be the prime components above $\Delta$, and
choose $N>0$ annihilating the finite group $H^2(Y;\Z)$.
The logarithms, in the exponential identification
$\Pic^0(Y)=H^1(Y,\OO_Y)$, of
$\OO_Y(ND_1),\ldots,\OO_Y(ND_r)$ generate a free abelian
subgroup of rank $r-s+1=t'$.
Here is the rank computation with its maps specified. The map is
$\Z^r\to\Pic^0(Y)$, $(n_i)\mapsto\OO_Y(N\sum n_iD_i)$.
A relation means that $N\sum n_iD_i$ is the divisor of a
meromorphic function. Since $f$ is the algebraic reduction, that
function is $f^*h$ for $h\in\mathcal M(\PP^1)$.
Its divisor is a sum of full fibres with total base degree zero.
Such relations have rank $s-1$, and those given by divisors of
rational functions supported on $\Delta$ already have this rank.
Statement \ref{cdp:st:D.1.2} gives the last equality.

\item Over $U=f^{-1}(C^\circ)$ each of these divisor line bundles has
its canonical trivialization. Exactness of the exponential sequence
therefore places its logarithm in the coefficient image of
$H^1(U;\Z)$. Coherent Leray identifies restriction
$H^1(Y,\OO_Y)\to H^1(U,\OO_U)$ with restriction of global
sections of $E$; this is injective because \CompendiumRef{II}{cdp:st:E4} makes $E$ locally free.
Thus the rank-$t'$ group of logarithms remains injective on $U$.

\item Topological Leray on the punctured curve gives
\[
0\to H^1(C^\circ;\Z)\to H^1(U;\Z)
\to H^1(F;\Z)^{\pi_1(C^\circ)}\to0.
\]
The first group's coefficient image vanishes because
$H^1(C^\circ,\OO)=0$. Conversely a class whose coefficient image
in $H^1(U,\OO)$ vanishes restricts to zero in $H^1(F,\OO_F)$;
Lemma \ref{cdp:st:D.1.7} makes its integral restriction to $F$ zero. Thus the
coefficient image of $H^1(U;\Z)$ identifies exactly with the
invariant lattice in \eqref{cdp:eq:D4}. The logarithm group maps to a subgroup
of the same rank $t'$, hence of finite index, in that lattice.
Its complex span equals the span of the entire lattice, and the
restriction image is a complex vector subspace. This proves the
assertion and the statement about the last spanning step.

\end{enumerate}
\end{proof}
